\documentclass[11pt,letterpaper]{article}
\usepackage[T1]{fontenc}
\usepackage[utf8]{inputenc}
\usepackage{lmodern}
\usepackage[margin=1in,headheight=14pt]{geometry}
\usepackage{amsmath,amssymb,amsthm,mathtools}
\usepackage{microtype,booktabs,longtable,array,enumitem}
\usepackage{xcolor,fancyhdr}
\usepackage[hyphens]{url}
\usepackage[colorlinks=true,linkcolor=blue!45!black,citecolor=blue!45!black,urlcolor=blue!45!black]{hyperref}
\hypersetup{pdftitle={Diophantine Laws of Probabilistic Computation},pdfauthor={ChatGPT, research manuscript prepared for Vladimir Reshetnikov},pdfsubject={Optimal normalization budgets, single-fold quartic certificates, and quantum computation}}
\pagestyle{fancy}
\fancyhf{}
\fancyhead[L]{\small Diophantine laws of probabilistic computation}
\fancyhead[R]{\small Research manuscript}
\fancyfoot[C]{\thepage}
\setlength{\parskip}{4pt}
\setlength{\parindent}{1.25em}
\setlength{\emergencystretch}{2em}
\allowdisplaybreaks[2]
\numberwithin{equation}{section}
\newtheorem{theorem}{Theorem}[section]
\newtheorem{lemma}[theorem]{Lemma}
\newtheorem{proposition}[theorem]{Proposition}
\newtheorem{corollary}[theorem]{Corollary}
\theoremstyle{definition}
\newtheorem{definition}[theorem]{Definition}
\newtheorem{example}[theorem]{Example}
\newtheorem{question}{Research question}
\theoremstyle{remark}
\newtheorem{remark}[theorem]{Remark}
\newcommand{\N}{\mathbb N}
\newcommand{\Z}{\mathbb Z}
\newcommand{\Q}{\mathbb Q}
\newcommand{\R}{\mathbb R}
\newcommand{\C}{\mathbb C}
\newcommand{\ce}{c.e.}
\newcommand{\TV}{\operatorname{TV}}
\newcommand{\Tr}{\operatorname{Tr}}
\newcommand{\supp}{\operatorname{supp}}
\newcommand{\diag}{\operatorname{diag}}
\newcommand{\Dmax}{D_{\max}}
\newcommand{\Dinf}{D_{\infty}}
\newcommand{\indic}[1]{\mathbf 1_{\{#1\}}}
\newcommand{\floor}[1]{\left\lfloor #1\right\rfloor}
\newcommand{\code}[1]{\texttt{\detokenize{#1}}}
\newcommand{\sigone}{\Sigma^0_1}
\newcommand{\pione}{\Pi^0_1}
\newcommand{\sigtwo}{\Sigma^0_2}
\newcommand{\pitwo}{\Pi^0_2}
\newcommand{\snapshot}{\texttt{f608f1cb3c5be8a736df1328c01b293aabfccf4e}}

\begin{document}
\begin{titlepage}
\centering
{\small PROVEIT-RELATED RESEARCH \quad / \quad 30 SEPTEMBER 2026\par}
\vspace{1.0cm}
{\LARGE\bfseries Diophantine Laws of\\[0.15cm] Probabilistic Computation\par}
\vspace{0.45cm}
{\large Optimal normalization budgets, unique quartic certificates,\\
and quantum computational substrates\par}
\vspace{0.6cm}
{\normalsize Prepared for Vladimir Reshetnikov\par}
{\small Mathematical development and reproducible implementation by ChatGPT\par}
\vspace{0.65cm}
\begin{minipage}{0.93\textwidth}
\begin{abstract}
The statement that the halting set of an effective computational substrate is
Diophantine follows from MRDP. For probabilistic and quantum computation, the
more discriminating questions concern output \emph{laws}, conditional
probabilities, and the cost of realizing prescribed normalized histories.
We develop a substrate-independent representation of an effective discrete
output law as the dyadic mass of a prefix-free projection of a fixed quartic
integer zero set. This construction counts projected events, not polynomial
witnesses, and therefore requires no finite-fold hypothesis.

The central constructive result determines the exact mass budget for a
sequence of strictly positive probability vectors. Its cost is the product
of the successive largest coordinate ratios. For a finite rational history
with $N$ transitions and $m$ outcomes, we construct one quartic polynomial
with $5Nm+m+4N+3$ natural witnesses. A valid feasible input has exactly one
witness tuple. A canonical earliest-maximum rule removes ties without
increasing the degree. We also prove a sharp binary variation bound and an
operator version governed by max-relative entropy.

Normalization changes the arithmetic boundary. Conditional two-outcome
probabilities with any fixed success floor below one realize exactly the
weakly computable reals in $[0,1]$, whereas almost-sure return realizes exactly
the computable reals. Strict conditional comparisons remain
$\Sigma^0_2$-complete even when success is guaranteed arbitrarily close to
one. Explicit integer arithmetic for Clifford+$T$ circuits supplies a quantum
front end. The accompanying code exports a concrete quartic and performs
exhaustive, symbolic, and exact-arithmetic tests. We distinguish these proved
constructions from established ingredients, unexpanded MRDP polynomials,
and formalization work that has not been performed.
\end{abstract}
\end{minipage}
\vfill
\begin{minipage}{0.93\textwidth}
\small
\textbf{Status.} This is a research manuscript with complete ordinary
mathematical proofs and executable finite checks, not a claim of independently
established historical priority or a Lean/Rocq-verified extension. The most
concrete deliverable is the explicit single-fold finite-history compiler in
Section~\ref{sec:quartic}. No general finite-fold MRDP conjecture is assumed or
settled.
\end{minipage}
\end{titlepage}
\tableofcontents
\newpage

\section{Research target and relation to ProveIt}\label{sec:context}

\subsection{What is being represented?}
A computational substrate may manipulate tape symbols, lambda terms,
fractions, rewrite graphs, amplitudes, or physical-looking states. If its
finite executions are effective, existential arithmetic can describe their
occurrence. This observation alone does not preserve the quantitative
semantics of probabilistic computation. A polynomial may have many witnesses
for one execution, and infinitely many distinct executions may contribute
to an output probability. Quantum computational paths can interfere, so
summing their squared amplitudes individually is generally wrong.

This article separates four questions. First, can one recover the \emph{whole
discrete output law} from integer polynomial solvability? Second, which
probability comparisons are themselves existential Diophantine? Third, what
is the exact probability mass required to realize a changing conditional
output distribution? Fourth, can this last optimization problem have a
small, explicit polynomial certificate with uniquely determined witnesses?
The answers form a coherent extension of the finite-trace viewpoint.

Our strongest concrete construction is not merely an application of MRDP.
It is a direct quartic compiler for optimal normalization feasibility,
including precise witness and equation counts, invalid-input rejection,
canonical tie resolution, and a uniqueness proof. The infinite and quantum
extensions explain why that finite compiler is useful.

\subsection{Repository boundary}
The inspected ProveIt snapshot is
\begin{center}\small\snapshot.\end{center}
Its Hilbert-tenth-problem development exposes the theorems
\code{mrdp}, \code{mrdp_iff}, and \code{mrdp_dioph_iff}
in the \code{Diophantine} namespace. In particular, the source of
\code{Diophantine/MRDP.lean} states the equivalence between recursively
enumerable subsets of $\N$ and projections of zero sets of integer
polynomials with a finite natural witness dimension, chosen before the
input. The statement includes input zero \cite{proveit-mrdp,proveit-readme}.
We use exactly that mathematical interface, not the correctness of every
other research claim in the repository.

The repository also supplies pairing, primitive-recursive graph, and trace
modules adjacent to this interface. Those are natural integration points
for the proposed compiler. This manuscript does not report a local rebuild
of ProveIt or an axiom audit of newly supplied declarations: it supplies no
new kernel-checked declarations. Its mathematical proofs are independent
of a choice of proof-assistant syntax.

\subsection{Established ingredients and contribution ledger}
The MRDP equivalence is an established theorem \cite{matiyasevich}. Formalized reductions through
FRACTRAN already exist \cite{larchey-forster}; hence a bare assertion that
FRACTRAN halting is Diophantine would not be a new result. Weakly computable
reals and their bounded-variation characterization also predate this work
\cite{ambos-spies,ng-stephan-wu}. Max-relative entropy and its operator-order
formula are established \cite{datta}. Finally, the arithmetic hardness of
probabilistic termination and expectation comparison has substantial prior
literature \cite{kaminski-katoen}.

\begin{center}
\small
\begin{tabular}{p{0.28\textwidth}p{0.62\textwidth}}
\toprule
Component & Status in this manuscript\\
\midrule
MRDP; prefix allocation; weak computability; max-relative entropy &
Established ingredients, stated with attribution and used with explicit proofs
of the needed consequences.\\[3pt]
Exact normalization budget &
Derived here as a finite/infinite optimization theorem, with a canonical
minimal realization.\\[3pt]
Unique quartic compiler &
Explicit construction and proof here; implemented and checked on finite
instances. Historical priority is not asserted.\\[3pt]
Sharp binary and operator bounds &
Proved here; the underlying entropy notions are not claimed as new.\\[3pt]
High-success arithmetic boundary &
Explicit realization and hardness reductions here, situated relative to
known computability classifications.\\[3pt]
Universal output-law polynomial &
Existence and degree reduction proved; its full coefficients and witness
count are not generated.\\
\bottomrule
\end{tabular}
\end{center}

The questions answered here are precisely formulated mathematical questions
about the proposed extension. They are not presented as previously published
open problems whose priority history has been exhaustively established.

\section{Effective output laws and quartic arithmetic}\label{sec:laws}

\subsection{Conventions}
We take $\N=\{0,1,2,\ldots\}$. A relation on natural tuples is
\emph{Diophantine} when it has the form
\begin{equation}\label{eq:dioph}
 R(\mathbf x)\quad\Longleftrightarrow\quad
 \exists\mathbf w\in\N^k\;P(\mathbf x,\mathbf w)=0,
 \qquad P\in\Z[\mathbf X,\mathbf W].
\end{equation}
All polynomial witnesses in this article are natural numbers. Negative
coefficients are allowed. This is not a statement about rational witnesses,
or about unrestricted integer witnesses unless an explicit translation is
given.

A real $a$ is \emph{left-c.e.} if a computable nondecreasing rational sequence
converges to $a$. A discrete \emph{subprobability law} is a function
$\mu:\N\to[0,1]$ with $\sum_z\mu(z)\le1$. Missing mass is allowed to represent
nonreturn. An effective family of output laws is uniformly left-c.e. in its
program and output arguments.

\begin{definition}[Effective first-return presentation]\label{def:firstreturn}
An effective first-return presentation consists of an effective enumeration
of finite terminal histories, a computable output label for each history,
and uniformly computable nonnegative real history weights. Histories
represent pairwise disjoint terminal outcomes, and the sum of their weights
is at most one. Each actual finite terminal history occurs once in the
presentation, or is canonically deduplicated.
\end{definition}

For a quantum machine, a history in this definition is a sequence of actual
measurement outcomes together with the classical control information.
Interfering basis paths between measurements must already have been combined
when computing its weight.

\begin{lemma}[Effective first returns give lower approximations]\label{lem:leftce}
Every effective first-return presentation has a uniformly left-c.e. output
law.
\end{lemma}
\begin{proof}
For each history $j$, compute rational lower bounds on its weight, replace
negative bounds by zero, and take running maxima. At stage $s$, sum the bounds
for the first $s$ enumerated histories with the specified output label.
These sums are computable, nondecreasing, and bounded by the true output
mass. Every finite collection of history weights is approximated arbitrarily
well, so the limit is their full sum. The construction is uniform. Rounding
down on successively finer dyadic grids gives dyadic lower approximations.
For infinitely many labels, expose only finitely many labels at each stage;
all lower masses still have sum at most one.
\end{proof}

The assumptions permit an infinite state space and an unbounded running
time. They do not permit noncomputable real transition constants hidden in
the source language. They also do not turn an infinite limiting configuration
into a finite first return.

\subsection{A degree-reduction interface}
\begin{lemma}[Quartic lowering]\label{lem:quarticlower}
Every Diophantine relation has a representation by a single integer
polynomial of total degree at most four. A fixed input polynomial gives a
fixed finite new witness dimension.
\end{lemma}
\begin{proof}
Evaluate the input polynomial by a finite arithmetic circuit. Introduce
variables for its intermediate values. Each addition has a linear residual,
and each multiplication has a quadratic residual. Signed intermediate
integers can be represented by differences of natural variables; a product
of two such differences is still quadratic. Equate the output difference
to zero and sum the squares of all residuals. An integer sum of squares is
zero exactly when every residual is zero. Conversely, a genuine evaluation
supplies intermediate values and hence natural positive/negative parts.
The resulting sum has total degree at most four. No uniqueness of the
introduced witnesses is required for this lemma.
\end{proof}

The polynomial in Lemma~\ref{lem:quarticlower} can be much larger than the
original one. Degree reduction is not an efficiency theorem. The direct
construction in Section~\ref{sec:quartic}, in contrast, controls its arithmetic
circuit size and natural witness multiplicity.

\section{Output laws as masses of Diophantine projections}\label{sec:prefix}

\subsection{Canonical allocation of fair-coin events}
A finite binary word $\sigma$ specifies the cylinder of infinite fair-coin
streams beginning with $\sigma$. Its probability is $2^{-|\sigma|}$.
A prefix-free set of words specifies disjoint cylinders.

\begin{lemma}[Effective labeled allocation]\label{lem:kraft}
Suppose a computable enumeration requests labeled binary words of lengths
$n_1,n_2,\ldots$, and
$\sum_j2^{-n_j}\le1$. There is a deterministic effective allocation of
prefix-free words with exactly these lengths and labels.
\end{lemma}
\begin{proof}
Maintain a finite set of free words, initially containing the empty word.
Maintain the invariant that there is at most one free word of each length.
For a request of length $n$, choose the longest free word of length at most
$n$. Remove it. Repeatedly split its cylinder into its left child, which is
continued, and its right child, which is put in the free set, until length
$n$ is reached. Allocate the final left descendant to the request.

There was no other free word with length strictly between the chosen length
and $n$, by maximality. Thus the one-free-word-per-length invariant is
preserved. Free and allocated cylinders remain pairwise disjoint and together
have total mass one. If there were no suitable free word, all free words
would have distinct lengths greater than $n$. Since the free set is finite,
its total mass would be strictly less than
$\sum_{k>n}2^{-k}=2^{-n}$. The requested mass would then exceed the remaining
mass, contrary to the hypothesis. Every allocation therefore succeeds.
\end{proof}

\begin{theorem}[Fair-coin realization]\label{thm:semimeasure}
A uniformly given discrete subprobability law is realizable as the terminal
output law of a fair-coin program if and only if it is uniformly left-c.e.
\end{theorem}
\begin{proof}
A fair-coin machine supplies an effective enumeration of disjoint finite
halting cylinders, so the forward implication follows directly, or from
Lemma~\ref{lem:leftce}.

For the reverse implication choose monotone dyadic lower approximations
$d_s(z)$, with finite support at each stage and total mass at most one.
Each increment $d_s(z)-d_{s-1}(z)$ is a nonnegative dyadic rational and can
be decomposed into finitely many requests for masses $2^{-n}$, all labeled
by $z$. The total requested mass is $\sum_z\mu(z)\le1$.
Lemma~\ref{lem:kraft} allocates their cylinders. A program dovetails this
allocation enumeration with reading its random stream and returns $z$ when
it discovers an allocated cylinder containing that stream. Its output events
are exactly the allocated cylinders, and their mass for label $z$ is
$\lim_s d_s(z)=\mu(z)$.
\end{proof}

The program may read additional random bits before an allocated cylinder is
recognized. The allocated word is an event code, not necessarily its literal
minimal runtime leaf. The theorem preserves the terminal output law, not
runtime, space, or interactive behavior.

\subsection{One quartic zero set for all effective laws}
Encode a binary word by
\begin{equation}
 c(\sigma)=2^{|\sigma|}+\operatorname{value}_2(\sigma),\qquad
 \ell(c)=\floor{\log_2 c}\quad(c\ge1).
\end{equation}
The empty word has code one. Fix a universal allocation interpreter. It
accepts requests only while preserving the prefix-free allocation invariant;
malformed or overbudget requests are discarded or leave the source stalled.
Thus arbitrary program codes still specify a valid subprobability law.
Valid source compilers from Theorem~\ref{thm:semimeasure} never trigger this
protection.

\begin{theorem}[Universal projected-root representation]\label{thm:universal}
There are a fixed $k$ and a fixed integer polynomial
$U(E,Z,C,W_1,\ldots,W_k)$ of total degree at most four such that every effective
discrete output law has, effectively from its presentation, an index $e$
with
\begin{equation}\label{eq:rootmass}
 \mu_e(z)=\sum_{c\ge1}2^{-\ell(c)}
 \indic{\exists\mathbf w\in\N^k\;U(e,z,c,\mathbf w)=0}.
\end{equation}
For fixed $e$, the word codes appearing in this sum across all labels are
prefix-free, and each code has at most one label.
\end{theorem}
\begin{proof}
The relation ``interpreter $e$ allocates word code $c$ to label $z$'' is c.e.:
run the interpreter and enumerate its allocations. Apply MRDP to this
relation, using an effective coding of its input tuple, and then apply
Lemma~\ref{lem:quarticlower}. The resulting polynomial and its witness
dimension depend only on the fixed interpreter. Prefix-freeness and label
uniqueness are properties of that interpreter. Equation~\eqref{eq:rootmass}
is the disjoint-cylinder probability formula. Every law from
Theorem~\ref{thm:semimeasure} has a source for the interpreter.
\end{proof}

\begin{remark}[Projection is essential]
Equation~\eqref{eq:rootmass} assigns one weight to a word code for which a
witness exists. It does \emph{not} sum over witness tuples. A single word may
have infinitely many arithmetic witnesses without changing its probability.
Consequently this theorem does not assume a finite-fold or single-fold MRDP
representation. Conversely, any effectively specified polynomial relation
whose output-labeled projection is prefix-free defines a left-c.e.
subprobability law by this same formula.
\end{remark}

\begin{corollary}[Ordinary strict thresholds]\label{cor:strict}
For any effective output-law interpreter, the relation
$b\mu_e(z)>a$, where $a\in\N$ and $b\ge1$, is Diophantine and has a fixed
quartic representation.
\end{corollary}
\begin{proof}
A monotone rational approximation exceeds $a/b$ at a finite stage exactly
when the limit does. This is a c.e. relation in $(e,z,a,b)$. Apply MRDP and
quartic lowering, including $b\ge1$ in the decidable input guard.
\end{proof}

Neither Theorem~\ref{thm:universal} nor Corollary~\ref{cor:strict} supplies the
full expanded universal polynomial in the accompanying software. They give
an existence construction through a fixed c.e. interpreter. The finite
polynomial exported by the software is the substantially more explicit one
constructed next.

\section{The exact normalization budget}\label{sec:budget}

\subsection{A mass-allocation problem}
Let $p_0,p_1,\ldots$ be strictly positive probability vectors on
$\{1,\ldots,m\}$, with $m\ge2$. Suppose that by checkpoint $s$ we have
accumulated masses $a_{s,i}$ of disjoint returned outcomes. Those masses must
be coordinatewise nondecreasing. Write
\begin{equation}\label{eq:normalization}
 h_s=\sum_i a_{s,i},\qquad a_{s,i}=h_sp_{s,i}.
\end{equation}
Here $h_s$ is the probability already assigned to a successful return, also
called the \emph{heralded mass}. Fix an initial mass $h_0=\delta>0$ and a
unit total budget $h_s\le1$.

The checkpoints are points in an effective enumeration of accumulated
probability. They need not coincide with successive machine instructions or
successive coin flips. This distinction is necessary: exact rational mass
increments need not be realizable in a uniformly bounded number of fair
coin flips.

\begin{definition}[Normalization cost]
For $s\ge1$, define
\begin{equation}\label{eq:R-G}
 R_s=\max_{1\le i\le m}\frac{p_{s-1,i}}{p_{s,i}},\qquad
 G_0=1,\qquad G_s=\prod_{j=1}^{s}R_j.
\end{equation}
The path's total multiplicative cost is $G_N$ for a finite horizon, and
$G_\infty=\sup_sG_s$ for an infinite history.
\end{definition}

Since both vectors sum to one, $R_s\ge1$. Otherwise every old coordinate
would be strictly smaller than the new coordinate, contradicting equal
sums.

\begin{theorem}[Optimal normalization budget]\label{thm:budget}
A finite or infinite positive probability history admits masses satisfying
\eqref{eq:normalization}, $h_0=\delta$, $h_s\le1$, and
$a_{s,i}\ge a_{s-1,i}$ if and only if
\begin{equation}\label{eq:budgetcondition}
 \delta G_s\le1\qquad\text{for every checkpoint }s.
\end{equation}
Its unique pointwise minimal feasible history is
\begin{equation}\label{eq:greedy}
 h_s^*=\delta G_s,\qquad a_{s,i}^*=\delta G_sp_{s,i}.
\end{equation}
In particular, the largest feasible initial mass for a finite path is
$G_N^{-1}$.
\end{theorem}
\begin{proof}
For a single transition, coordinatewise monotonicity is equivalent to
\[
 h_sp_{s,i}\ge h_{s-1}p_{s-1,i}\quad\text{for every }i,
\]
and, because all quantities being divided by are positive, to
$h_s\ge h_{s-1}R_s$. Iterating gives $h_s\ge\delta G_s$.
This proves necessity and pointwise minimality of any possible candidate.

Conversely, the choice in \eqref{eq:greedy} satisfies the single-step
inequality with equality in the maximizing ratio and inequality in all
other coordinates. Its coordinates are nondecreasing, its initial mass is
$\delta$, and the assumed budget condition gives total mass at most one.
It is therefore feasible. Each of its total masses is forced for any
pointwise minimal solution, and the prescribed vectors then force every
coordinate. The argument applies to every finite prefix of an infinite
history without any compactness or choice step.
\end{proof}

\begin{corollary}[Effective realization of a feasible history]\label{cor:effectivebudget}
If $\delta$ and the $p_{s,i}$ are computable rationals, and the budget
condition holds, the greedy masses define a computable nondecreasing
rational approximation to an output subprobability law. A fair-coin program
realizes that limiting law.
\end{corollary}
\begin{proof}
Finite maxima and products of positive rationals are computable exactly.
Theorem~\ref{thm:budget} supplies monotonicity and boundedness. Each coordinate
therefore converges and is left-c.e.; the limiting sum is at most one.
Apply Theorem~\ref{thm:semimeasure}. When necessary, round rational cumulative
masses down to finer dyadic lower approximations before allocation.
\end{proof}

\begin{example}[An excursion costs mass]\label{ex:excursion}
Consider
\[
 p_0=(1/2,1/2),\quad p_1=(2/3,1/3),\quad p_2=(1/2,1/2).
\]
Then $R_1=3/2$, $R_2=4/3$, and $G_2=2$. Thus $\delta\le1/2$ is necessary
and sufficient, although the initial and final conditional laws agree.
At the maximum initial mass, the greedy cumulative masses are
\[
 (1/4,1/4),\quad(1/2,1/4),\quad(1/2,1/2).
\]
The first transition spends additional mass on the first output; the return
transition spends it on the second. Previously assigned mass cannot be
withdrawn.
\end{example}

\subsection{An information-theoretic expression}
For positive classical probability vectors, write
\[
 \Dinf(p\Vert q)=\log\max_i\frac{p_i}{q_i},
\]
with the natural logarithm throughout this article. Then
\begin{equation}\label{eq:costentropy}
 \log G_N=\sum_{s=1}^{N}\Dinf(p_{s-1}\Vert p_s).
\end{equation}
The divergence itself is standard, and is the commuting specialization of
max-relative entropy \cite{datta}. Its role here is an exact additive budget
for irreversible accumulation of terminal mass. The direction of the
arguments is significant: the old law is compared with the new law.
Reversing them generally changes the answer.

\section{A single-fold quartic certificate with explicit size}\label{sec:quartic}

\subsection{Parameters and statement}
Fix an external horizon $N\ge0$ and an external number of labels $m\ge2$.
The free natural parameters are
\[
 (u_{s,i})_{0\le s\le N,\,1\le i\le m},\quad
 (v_s)_{0\le s\le N},\quad \alpha,\beta.
\]
Call an input \emph{valid} when
\begin{equation}\label{eq:valid}
 u_{s,i}\ge1,\quad v_s=\sum_i u_{s,i},\qquad
 1\le\alpha\le\beta.
\end{equation}
It specifies $p_{s,i}=u_{s,i}/v_s$ and $\delta=\alpha/\beta$.
Denominators need not be reduced, and different rows may use independently
rescaled denominators.

\begin{theorem}[Explicit unique-witness quartic]\label{thm:singlefold}
For each fixed $N,m$ there is an explicitly constructible integer polynomial
$F_{N,m}$ of total degree at most four with
\begin{align}
 W(N,m)&=5Nm+m+4N+3 \label{eq:wcount}\\
 E(N,m)&=5Nm+m+6N+4 \label{eq:ecount}
\end{align}
natural witnesses and squared residuals, respectively, such that:
\begin{enumerate}[label=(\roman*),nosep]
\item a valid input has a witness if and only if its prescribed history is
feasible with initial mass $\alpha/\beta$;
\item every valid feasible input has exactly one natural witness tuple;
\item invalid inputs have no natural witness tuple.
\end{enumerate}
The residual system has an arithmetic-circuit implementation of size
$O(Nm)$. This is not a claim that its fully expanded polynomial has only
$O(Nm)$ monomials.
\end{theorem}

\subsection{Validity, ratios, and selectors}
Introduce positivity witnesses $g_{s,i}$ and $g_\alpha,g_\beta$, with residuals
\begin{align}
 u_{s,i}-1-g_{s,i}&=0,&v_s-\sum_i u_{s,i}&=0,\label{eq:validres}\\
 \alpha-1-g_\alpha&=0,&\beta-\alpha-g_\beta&=0.\notag
\end{align}
For every transition $s=1,\ldots,N$ introduce $X_{s,i},Y_{s,i}$ and enforce
\begin{equation}\label{eq:XY}
 X_{s,i}-u_{s-1,i}v_s=0,\qquad
 Y_{s,i}-v_{s-1}u_{s,i}=0.
\end{equation}
Their ratios are the candidates for $R_s$:
$X_{s,i}/Y_{s,i}=p_{s-1,i}/p_{s,i}$, and their denominators are positive.

Introduce selectors $e_{s,i}$ satisfying
\begin{equation}\label{eq:selector}
 e_{s,i}(e_{s,i}-1)=0,\qquad \sum_i e_{s,i}-1=0.
\end{equation}
They select exactly one coordinate. Define selected numerator and denominator
witnesses $A_s,B_s$ by
\begin{equation}\label{eq:AB}
 A_s-\sum_i e_{s,i}X_{s,i}=0,\qquad
 B_s-\sum_i e_{s,i}Y_{s,i}=0.
\end{equation}
Every residual displayed so far has total degree at most two.

\subsection{Canonical earliest-maximum rule}
Ordinary maximum constraints allow several selectors when two ratios tie.
To eliminate this multiplicity, introduce natural slacks $t_{s,i}$ and use
\begin{equation}\label{eq:canonical}
 A_sY_{s,i}-B_sX_{s,i}-t_{s,i}-\sum_{j>i}e_{s,j}=0.
\end{equation}
If coordinate $k$ is selected, the last sum is one for $i<k$ and zero for
$i\ge k$. Thus the selected ratio must be \emph{strictly} larger than every
earlier ratio and at least as large as every later one. Integer strictness
is encoded by subtraction of one, not by a new quantifier or a high-degree
product.

\begin{lemma}[Unique pivot]\label{lem:pivot}
For any positive integer lists $(X_i),(Y_i)$, the selector and comparison
constraints \eqref{eq:selector}--\eqref{eq:canonical} have exactly one solution
in the selector, selected-ratio, and slack variables. Its selected index is
the least index maximizing $X_i/Y_i$.
\end{lemma}
\begin{proof}
The Boolean and sum constraints select an index $k$. Nonnegative slacks give
$X_kY_i-Y_kX_i\ge1$ for $i<k$ and $\ge0$ for $i\ge k$.
Positivity of all denominators makes these exactly the asserted strict and
weak ratio comparisons. Hence $k$ must be the least maximizing index.
Conversely, that index satisfies those comparisons: a strictly positive
integer cross difference is at least one. It determines all selectors,
$A=X_k$, $B=Y_k$, and all slacks uniquely.
\end{proof}

\subsection{Canonical products and the terminal budget}
Set $C_0=\alpha$, $D_0=\beta$ as expressions, not additional witnesses.
For $s=1,\ldots,N$ introduce $C_s,D_s$ and impose
\begin{equation}\label{eq:CD}
 C_s-C_{s-1}A_s=0,\qquad D_s-D_{s-1}B_s=0.
\end{equation}
Finally introduce a natural slack $z$ and impose
\begin{equation}\label{eq:terminal}
 D_N-C_N-z=0.
\end{equation}
For $N=0$, this uses $D_0=\beta,C_0=\alpha$ directly. Products deliberately
remain unreduced. There are no arbitrary common factors and no gcd
certificates to create extra witnesses.

Let $\mathcal R_{N,m}$ be the complete residual list in
\eqref{eq:validres}--\eqref{eq:terminal}, including every indicated index, and
put
\begin{equation}\label{eq:F}
 F_{N,m}=\sum_{r\in\mathcal R_{N,m}}r^2.
\end{equation}
This is the promised explicit polynomial definition. It is a finite
sum-of-squares specification, not an appeal to MRDP.

\subsection{Proof of the compiler theorem}
\begin{proof}[Proof of Theorem~\ref{thm:singlefold}]
\emph{Soundness.} A natural zero of $F_{N,m}$ makes each residual zero.
The validity residuals force \eqref{eq:valid}. Equations~\eqref{eq:XY}
then force positive ratio denominators. Lemma~\ref{lem:pivot} forces the
least maximizing ratio at each stage, so $A_s/B_s=R_s$.
The product residuals force
\[
 \frac{C_N}{D_N}=\frac{\alpha}{\beta}\prod_{s=1}^N R_s
 =\delta G_N.
\]
The denominator is positive. The final natural slack forces
$C_N\le D_N$, hence $\delta G_N\le1$. Since $G_s$ is nondecreasing, the
budget condition holds for every earlier stage. Apply
Theorem~\ref{thm:budget}.

\emph{Completeness.} For a valid feasible input, assign the positivity and
input slacks their forced values, compute the $X,Y$ lists, choose each
least maximizing coordinate, and assign its uniquely determined nonnegative
comparison slacks. Compute the unreduced products $C_s,D_s$. Feasibility
makes $D_N-C_N$ nonnegative; assign it to $z$. Every residual is zero.

\emph{Uniqueness.} The input fixes the validity slacks and all $X,Y$ values.
Lemma~\ref{lem:pivot} fixes the selectors, $A,B$, and comparison slacks.
The recurrence fixes the unreduced products, and the final equation fixes
$z$. There are no other witnesses. The existence of other, nonminimal
feasible mass histories is irrelevant: the certificate describes the
canonical optimal budget, not an arbitrary feasible history.

\emph{Size and degree.} The positivity slacks contribute $(N+1)m$ witnesses,
the two input slacks contribute two, and each transition contributes
$4m+4$ witnesses: $X,Y,e,t$ for every coordinate and $A,B,C,D$.
The terminal slack contributes one. This gives \eqref{eq:wcount}.
Validity contributes $(N+1)(m+1)+2$ residuals, each transition contributes
$4m+5$, and the terminal equation contributes one, giving
\eqref{eq:ecount}. Each residual is quadratic or linear, so their squares
have degree at most four. Shared suffix sums of the selectors evaluate all
comparison residuals with $O(m)$ arithmetic operations per transition.
\end{proof}

\begin{remark}[The exact single-fold scope]
The horizon $N$ and label count $m$ index a family of finite polynomials.
They are not unbounded runtime variables inside one fixed-arity
single-fold universal halting polynomial. This construction therefore does
not resolve a general finite-fold MRDP question. Its uniqueness is over
natural witnesses. Replacing each natural variable by an arbitrary
four-square representation would in general destroy that uniqueness.
\end{remark}

\subsection{The exported \texorpdfstring{$N=2,m=2$}{N=2, m=2} polynomial}
For Example~\ref{ex:excursion}, take rows
$(u_{0,1},u_{0,2})=(1,1)$, $(u_{1,1},u_{1,2})=(2,1)$,
$(u_{2,1},u_{2,2})=(1,1)$ and $(\alpha,\beta)=(1,2)$.
The denominators are $(2,3,2)$. The selected coordinate is two at the first
transition and one at the second. Thus
\[
 (A_1,B_1)=(3,2),\quad(A_2,B_2)=(4,3),\qquad
 (C_1,D_1)=(3,4),\quad(C_2,D_2)=(12,12).
\]
The terminal slack is zero. This instance has 11 free input coordinates,
33 witnesses, 38 residuals, total degree four, and 165 expanded monomials.
The bundle supplies the residual list, variable ordering, full natural
assignment, and expanded polynomial. These values are computed by the
included script rather than inferred from a numerical optimizer.

\section{Variation bounds and a sharp binary constant}\label{sec:variation}

\subsection{A general bound}
For probability vectors define
$\TV(p,q)=\tfrac12\sum_i|p_i-q_i|$.

\begin{proposition}[Variation consumes mass]\label{prop:variation}
For the positive histories of Theorem~\ref{thm:budget},
\begin{equation}\label{eq:tvbound}
 \sum_{s=1}^{N}\TV(p_{s-1},p_s)
 \le\sum_{s=1}^{N}(1-R_s^{-1})
 \le\log G_N.
\end{equation}
Every history feasible from initial mass $\delta$ therefore has total
variation at most $\log(1/\delta)$, including an infinite history.
\end{proposition}
\begin{proof}
If $R_s=1$, then $p_s\ge p_{s-1}$ coordinatewise and equal total sums imply
equality. If $R_s>1$, coordinatewise domination gives a probability vector
$q_s$ such that
\[
 p_s=R_s^{-1}p_{s-1}+(1-R_s^{-1})q_s.
\]
Hence $\TV(p_{s-1},p_s)\le1-R_s^{-1}$. For $R\ge1$,
$1-R^{-1}\le\log R$, as follows by integrating $1/t$ on $[1,R]$ or
differentiating their difference. Sum over $s$ and apply the budget theorem.
For an infinite path pass to the supremum of the nonnegative partial sums.
\end{proof}

\subsection{The binary path identity}
The binary case contains a stronger symmetric bound. Put
$p_s=(x_s,1-x_s)$ with $0<x_s<1$, and define
\begin{equation}\label{eq:Phi}
 \Phi(x)=-\frac12\log\bigl(x(1-x)\bigr).
\end{equation}

\begin{theorem}[Sharp balanced binary bound]\label{thm:binary}
Every binary path satisfies
\begin{equation}\label{eq:binarypotential}
 2\sum_{s=1}^{N}|x_s-x_{s-1}|+
 \Phi(x_N)-\Phi(x_0)\le\log G_N.
\end{equation}
In particular, a feasible path starting at $x_0=1/2$ and initial mass
$\delta$ satisfies
\begin{equation}\label{eq:binarysharp}
 \sum_{s=1}^{N}|x_s-x_{s-1}|
 \le\frac12\log(1/\delta).
\end{equation}
For every $0<\delta<1$, the constant $1/2$ in this statement is sharp as a
supremum over finite balanced-start paths.
\end{theorem}
\begin{proof}
If a step increases $x$ from $x$ to $y$, its largest ratio is
$(1-x)/(1-y)$, so its logarithmic cost is
$\int_x^y(1-u)^{-1}\,du$. If the step decreases $x$ to $y$, the cost is
$\int_y^x u^{-1}\,du$. Since
\[
 \Phi'(u)=\frac{2u-1}{2u(1-u)},
\]
both cases have the exact common expression
\begin{equation}\label{eq:pathidentity}
 \log R_s=
 \int_{\min(x_{s-1},x_s)}^{\max(x_{s-1},x_s)}
 \frac{du}{2u(1-u)}+\Phi(x_s)-\Phi(x_{s-1}).
\end{equation}
Now $u(1-u)\le1/4$, so the integral is at least
$2|x_s-x_{s-1}|$. Sum and telescope the potential to obtain
\eqref{eq:binarypotential}. The potential is minimized at $1/2$, and
$\log G_N\le\log(1/\delta)$ for a feasible path, proving
\eqref{eq:binarysharp}.

For sharpness, repeat the cycle $1/2\to1/2+\varepsilon\to1/2$ exactly $n$
times. Its one-cycle multiplicative cost is
$(1+2\varepsilon)/(1-2\varepsilon)$ and its variation is $2\varepsilon$.
Choose
\[
 r_n=(1/\delta)^{1/n},\qquad
 \varepsilon_n=\frac{r_n-1}{2(r_n+1)}.
\]
The full cost is exactly $1/\delta$, while the full variation is
$n(r_n-1)/(r_n+1)$, which tends to
$\tfrac12\log(1/\delta)$. Thus no smaller constant works. If rational paths
are required, approximate each $\varepsilon_n$ from below by rationals;
feasibility is preserved and the same supremum is obtained.
\end{proof}

This bound is a quantitative restriction on successive conditional
estimates. Conditioning can hide noncomputable limiting information, but
it cannot make a fixed positive amount of accumulated mass support an
arbitrarily oscillatory normalized path. Small excursions near balance are
the asymptotically least expensive way to accumulate binary variation.

\section{Which real probabilities can normalization produce?}\label{sec:weak}

\subsection{Weak computability and finite variation}
A real is \emph{weakly computable}, also called a difference of left-c.e.
reals, when $c=A-B$ for finite left-c.e. reals $A,B$. The equivalence with
computable rational approximation of bounded total variation is established
in the literature \cite{ambos-spies,ng-stephan-wu}. The elementary argument
needed here is short.

\begin{lemma}[Jordan decomposition of an approximation]\label{lem:jordan}
A real $c$ is weakly computable if and only if there is a computable rational
sequence $(x_s)$ tending to $c$ for which
$\sum_s|x_{s+1}-x_s|<\infty$.
\end{lemma}
\begin{proof}
If $x_s=A_s-B_s$ with monotone convergent computable rational approximations,
then
$|x_{s+1}-x_s|\le(A_{s+1}-A_s)+(B_{s+1}-B_s)$, and the sum is finite.
Conversely, set
\[
 P_s=\sum_{j<s}\max(x_{j+1}-x_j,0),\qquad
 Q_s=\sum_{j<s}\max(x_j-x_{j+1},0).
\]
Both are computable increasing rational sequences with finite limits, and
$x_s=x_0+P_s-Q_s$. Their limits express $c$ as a difference of left-c.e.
reals.
\end{proof}

\subsection{An exact success-floor classification}
Two disjoint successful terminal outcomes, $A$ and $B$, have unconditional
masses $a,b$ and heralding probability $h=a+b$. Their conditional probability
of $A$, given eventual successful return, is
\begin{equation}\label{eq:conditional}
 c=\frac{a}{a+b},\qquad h>0.
\end{equation}
The denominator is an eventual-event probability. It is not assumed known
or computable.

\begin{theorem}[High-success normalization realizes exactly weak computability]
\label{thm:weakclassification}
For a real $c\in[0,1]$, the following are equivalent:
\begin{enumerate}[label=(\roman*),nosep]
\item $c$ is a conditional terminal-output probability of an effective
process with positive heralding probability;
\item $c$ is weakly computable;
\item for every rational $0<\delta<1$, there is such a process with heralding
probability at least $\delta$ and conditional probability exactly $c$.
\end{enumerate}
The realization in (iii) may be chosen to have a computable rational
accumulation checkpoint of total mass exactly $\delta$.
\end{theorem}
\begin{proof}
\emph{(i) implies (ii).} The masses $a,b$ are left-c.e. by
Lemma~\ref{lem:leftce}. If one is zero, $c$ is an endpoint and is computable.
Otherwise choose monotone rational approximations $a_s,b_s$ and discard a
finite initial segment until both are positive. Set
$h_s=a_s+b_s$ and $x_s=a_s/h_s$. These give an effective positive binary
history with initial mass $h_0>0$ and total budget at most one.
Proposition~\ref{prop:variation} bounds its total variation by
$\log(1/h_0)$. It converges to $a/(a+b)$, so
Lemma~\ref{lem:jordan} proves weak computability.

\emph{(ii) implies (iii).} The endpoint cases are realized by a constant
output. Suppose $0<c<1$, and take a computable rational bounded-variation
approximation to $c$. Choose rational $\eta>0$ with
$c\in(\eta,1-\eta)$. Clamp the approximation to $[\eta,1-\eta]$; clamping
is a contraction, preserves the limit, and does not increase total variation.
Discard a sufficiently long finite initial segment so that the remaining
variation is at most $\eta\log(1/\delta)$.

For a binary step confined to that interval, the integral formulas in the
proof of Theorem~\ref{thm:binary} give
$\log R_s\le|x_s-x_{s-1}|/\eta$. Hence the remaining path has
$\sum_s\log R_s\le\log(1/\delta)$. Its greedy realization from initial
mass $\delta$ is therefore feasible by Theorem~\ref{thm:budget}. Its
monotone rational coordinate masses converge to $a,b$ with $a+b\ge\delta$.
Because their normalized ratios are the chosen $x_s$, continuity at the
positive denominator gives $a/(a+b)=c$. Apply
Corollary~\ref{cor:effectivebudget} to realize the limiting law.

\emph{(iii) implies (i)} is immediate.
\end{proof}

\begin{remark}[The nonuniform choice must not be hidden]\label{rem:nonuniform}
The proof of (ii)$\Rightarrow$(iii) chooses an interior interval and a late
small-variation tail. It is an existence theorem for each real and each
success floor; it is not an algorithm for discovering those choices from
an arbitrary bounded-variation approximation index. The rational greedy
compiler is uniform once a suitable tail and a valid budget bound are
supplied. The finite-horizon quartic compiler is fully explicit and has no
such unbounded tail choice.
\end{remark}

Weakly computable reals are not all limit-computable reals; the distinction
is established by computability-theoretic results such as those discussed
by Ng, Stephan, and Wu \cite{ng-stephan-wu}. Thus the theorem does not say
that conditioning can realize an arbitrary $\Delta^0_2$ real. It identifies
an exact intermediate class using a concrete mass constraint.

\begin{theorem}[The full-success boundary]\label{thm:total}
The conditional probabilities of effective two-outcome processes that
return one of the two outcomes with probability one are exactly the
computable reals in $[0,1]$.
\end{theorem}
\begin{proof}
Here $a+b=1$. Increasing rational lower approximations give intervals
$[a_s,1-b_s]$ containing $a$, whose widths tend to zero. To compute $a$ to
any required precision, search for an interval of sufficiently small width.
The search terminates under the probability-one promise. Conversely, for a
computable $c$, both $c$ and $1-c$ are uniformly left-c.e. Their two-label
law has total mass one and is realized by
Theorem~\ref{thm:semimeasure}.
\end{proof}

There is therefore a genuine endpoint distinction: every fixed permitted
failure probability greater than zero allows all weakly computable
conditional probabilities; requiring failure probability exactly zero
restricts the class to computable probabilities. This is an expressivity
statement, not a physical procedure for obtaining a noncomputable real to
certified arbitrary precision from experimental samples.

\section{The arithmetic hierarchy of probability questions}\label{sec:hierarchy}

\subsection{Indexing and promises}
A $\sigone$ relation has an existential quantifier over a computable
predicate; $\sigtwo$ has an existential block followed by a universal block.
The dual classes are $\pione$ and $\pitwo$. Completeness below means
computable many-one completeness. Inputs are program indices for a fixed
effective universal model, with rational threshold $q\in(0,1)$ fixed in
advance.

When a success floor is required, use an effective family of wrappers that
reserves a positive rational mass for guaranteed return before handing the
remaining mass to an arbitrary source program. Restrict the syntax to
wrappers whose reserved mass is at least the specified $\delta$.
All indices in that family satisfy the floor; no undecidable semantic
validation is included in the language being classified. Equivalently,
one may state the upper bounds as promise-problem bounds and ensure that
every hardness reduction lands inside the promise. We will do both
explicitly rather than silently conjoin an arbitrary-code problem with an
additional floor predicate.

These distinctions matter because almost-sure termination itself has high
arithmetic complexity, as is known from prior work
\cite{kaminski-katoen}. Our proofs below establish the particular threshold
claims directly.

\subsection{Unconditional probabilities}
\begin{theorem}[Unconditional comparison table]\label{thm:unconditionaltable}
For a fixed interior rational $q$ and a universal effective output-law model,
the following classifications hold:
\begin{center}
\begin{tabular}{cccccc}
\toprule
$p>q$ & $p\le q$ & $p<q$ & $p\ge q$ & $p=q$ & $p\ne q$\\
\midrule
$\sigone$-c. & $\pione$-c. & $\sigtwo$-c. & $\pitwo$-c. & $\pitwo$-c. & $\sigtwo$-c.\\
\bottomrule
\end{tabular}
\end{center}
Here ``-c.'' means complete for the displayed class.
\end{theorem}
\begin{proof}
Let $p_s\uparrow p$ be computable rational lower bounds. Then
$p>q\iff\exists s\,(p_s>q)$ and $p\le q$ is its complement. Moreover,
\[
 p\ge q\iff\forall k\;\exists s\;(p_s>q-2^{-k}),
\]
so $p\ge q$ is $\pitwo$ and $p<q$ is $\sigtwo$. Equality is the conjunction
of $p\ge q$ and $p\le q$, hence is $\pitwo$; inequality is its complement.

For the first hardness claim, use a rational baseline $q$ and add a small
rational mass $\varepsilon>0$ precisely when a given ordinary computation
halts. The resulting mass exceeds $q$ exactly in the halting case.
Taking complements proves the $\pione$ claim.

For the remaining claims, let $\varphi_e$ be the $e$th ordinary partial
computable function. A fair-coin program chooses $n\ge0$ with probability
$2^{-(n+1)}$, simulates $\varphi_e(n)$, and returns precisely if that
simulation halts. Its return probability is
\begin{equation}\label{eq:te}
 t_e=\sum_{n\ge0}2^{-(n+1)}\indic{\varphi_e(n)\downarrow}.
\end{equation}
Thus $t_e=1$ exactly when $e$ belongs to the $\pitwo$-complete totality set.
Set $p_e=qt_e$. Then $p_e<q$ and $p_e\ne q$ exactly for nontotal indices,
whereas $p_e\ge q$ and $p_e=q$ exactly for total indices. These are the
required reductions.
\end{proof}

For completeness, totality has the form
$\forall n\exists t$ ``$\varphi_e(n)$ halts by $t$'', and every
$\pitwo$ predicate $\forall n\exists t\,R(e,n,t)$ with computable $R$
reduces to it by the program that, on $n$, searches for such a $t$.
No oracle is used by the probabilistic program in \eqref{eq:te}.

\subsection{Conditioning defeats the existential threshold boundary}
\begin{theorem}[Conditional comparison at any subunit success floor]
\label{thm:conditionaltable}
Fix rationals $q\in(0,1)$ and $\delta\in(0,1)$. In a universal effective
wrapper family guaranteeing $a+b\ge\delta$, set $c=a/(a+b)$. Then
\begin{center}
\begin{tabular}{cccccc}
\toprule
$c>q$ & $c<q$ & $c\ne q$ & $c\ge q$ & $c\le q$ & $c=q$\\
\midrule
$\sigtwo$-c. & $\sigtwo$-c. & $\sigtwo$-c. & $\pitwo$-c. & $\pitwo$-c. & $\pitwo$-c.\\
\bottomrule
\end{tabular}
\end{center}
In particular, the strict conditional comparisons are not existential
Diophantine relations in such a universal family.
\end{theorem}
\begin{proof}
Choose computable rational lower approximations $a_s,b_s$. Their normalized
ratios $c_s=a_s/(a_s+b_s)$ are computable once the denominator is positive;
use any fixed rational value before that point. They converge to $c$.
Consequently
\begin{equation}\label{eq:strictsigma2}
 c>q\iff\exists k\exists n\;\forall s\ge n\quad c_s\ge q+2^{-k}.
\end{equation}
The analogous expression handles $c<q$. Thus these strict cuts are
$\sigtwo$, and their complements are $\pitwo$. Equality has the more useful
form
\begin{equation}\label{eq:eqpi2}
 c=q\iff\forall k\forall n\;\exists s\ge n\quad |c_s-q|<2^{-k}.
\end{equation}
Convergence makes repeated arbitrary closeness equivalent to equality, so
this is a $\pitwo$ expression. Its complement is $\sigtwo$.

For hardness choose a positive rational
$\varepsilon<\min(1-\delta,q,1-q)$ and use $t_e$ from \eqref{eq:te}.
The first wrapper has masses
\begin{equation}\label{eq:hardplus}
 a_e=q,\qquad b_e=1-q-\varepsilon+\varepsilon t_e.
\end{equation}
It reserves mass $1-\varepsilon$ for guaranteed return and assigns the
remaining mass to the totality sampler. Its heralding probability is
$h_e=1-\varepsilon+\varepsilon t_e\ge\delta$, and
$c_e=q/h_e\ge q$. It follows that $c_e>q$ and $c_e\ne q$ hold exactly
when $e$ is nontotal, while $c_e\le q$ and $c_e=q$ hold exactly when $e$
is total.

For the other direction use
\begin{equation}\label{eq:hardminus}
 a_e=q-\varepsilon+\varepsilon t_e,\qquad b_e=1-q.
\end{equation}
The same denominator gives
$c_e-q=-(1-q)\varepsilon(1-t_e)/h_e$. Thus $c_e<q$ exactly for nontotal
indices and $c_e\ge q$ exactly for total indices. All masses in both
wrappers are nonnegative and their sums are at most one. These reductions
are effective, preserve the syntactic floor guarantee, and prove every
hardness assertion. MRDP and the strict arithmetic hierarchy then exclude
ordinary existential Diophantine representations for the complete
$\sigtwo$ cuts.
\end{proof}

For example, the theorem applies with a success guarantee of $0.999999$.
The difficulty is not confined to rare postselection events: the unknown
remaining denominator mass can encode a universal nontermination question.
A finite lower bound for $a$ does not certify that future additions to $b$
will not reverse a conditional threshold comparison.

\begin{remark}[A quantifier trap]
The naive formula saying that $c_s$ is eventually within every tolerance
of $q$ has the pattern $\forall k\exists n\forall s$ and only gives a
$\Pi^0_3$ upper bound. Formula~\eqref{eq:eqpi2} is sharper because convergence
is already known. Omitting that argument would obscure the exact boundary.
\end{remark}

\subsection{The probability-one promise}
\begin{proposition}[Comparisons under almost-sure return]\label{prop:totalcuts}
Under the promise $a+b=1$, strict comparisons $c>q,c<q$ and inequality
$c\ne q$ are $\sigone$-complete promise problems. Weak comparisons and
equality are $\pione$-complete promise problems. The reductions can be
chosen from a uniformly total family of samplers.
\end{proposition}
\begin{proof}
Under the promise $c=a$ and $1-c=b$, so $c>q$ and $c<q$ are witnessed by
lower approximations to $a$ and $b$, respectively. Their disjunction is
inequality; taking complements within the promise gives the upper bounds.

Let $x_e=2^{-t}$ if a designated computation halts at time $t$, and let
$x_e=0$ if it never halts. This is a uniformly computable real: after $s$
steps, a still-undetected later halting time can contribute at most
$2^{-s}$. For a sufficiently small positive rational $\varepsilon$, the
probabilities $q+\varepsilon x_e$ and $q-\varepsilon x_e$ lie in $[0,1]$
and are uniformly computable. Their complementary two-label laws have
uniform total samplers by Theorem~\ref{thm:semimeasure}. Strict comparison
in the appropriate direction, or inequality with $q$, detects halting;
equality and the opposite weak comparison detect nonhalting.
\end{proof}

This is a promise classification. Deciding which arbitrary probabilistic
programs satisfy the promise has not become easy. It is incorrect to erase
that promise and infer a $\pione$ algorithm for arbitrary exact-output
questions.

\subsection{Polynomial quantifier normal forms}
The hierarchy can itself be expressed using quartic integer equations.
For a computable relation $R(e,u,v)$, MRDP applies both to $R$ and to its
complement. Consequently the four normal forms are
\begin{center}
\begin{tabular}{cl}
\toprule
Class & Quartic arithmetic form, with natural quantified tuples\\
\midrule
$\sigone$ & $\exists\mathbf w\;P(e,\mathbf w)=0$\\
$\pione$ & $\forall\mathbf w\;P(e,\mathbf w)\ne0$\\
$\pitwo$ & $\forall u\exists v,\mathbf w\;P(e,u,v,\mathbf w)=0$\\
$\sigtwo$ & $\exists u\forall v,\mathbf w\;Q(e,u,v,\mathbf w)\ne0$\\
\bottomrule
\end{tabular}
\end{center}
For the last line choose $Q$ to represent $\neg R$, so that universal
nonvanishing expresses $\forall v\,R(e,u,v)$. Apply
Lemma~\ref{lem:quarticlower} to each matrix. These are arithmetic formulas
\emph{over} Diophantine relations, not all existential Diophantine
representations. The strict hierarchy is exactly why the distinction
cannot generally be removed.

All tables assumed $0<q<1$. At endpoints some classifications collapse:
$p>0$ is $\sigone$, $p=0$ is $\pione$, $p=1$ can be $\pitwo$-complete,
and $p>1$ is impossible. Endpoint statements must not be read off the
interior tables without checking them.

\section{An explicit quantum computational front end}\label{sec:quantum}

\subsection{Exact amplitudes rather than floating-point paths}
The Clifford+$T$ gate family admits exact arithmetic in
$\Z[1/\sqrt2,i]$; the exact synthesis literature gives a much stronger
structural account of this ring and its realizable operators
\cite{giles-selinger}. Here we only need a direct forward evaluator for the
gates $H,T,S,X,$ and CNOT.

After $L$ elementary gates, represent each amplitude as
\begin{equation}\label{eq:ring}
 \psi_j=\frac{a_j+b_j\sqrt2+i(c_j+d_j\sqrt2)}{2^L},
 \qquad a_j,b_j,c_j,d_j\in\Z.
\end{equation}
The common denominator is intentionally unreduced: each gate increments
$L$ by one, even when it could be implemented without a new denominator.
For $z=a+b\sqrt2+i(c+d\sqrt2)$ the relevant integer-linear maps are
\begin{align}
 \sqrt2 z&\longleftrightarrow(2b,a,2d,c),\label{eq:linearquantum}\\
 iz&\longleftrightarrow(-c,-d,a,b),\notag\\
 \sqrt2(1+i)z&\longleftrightarrow
 \bigl(2(b-d),a-c,2(b+d),a+c\bigr).\notag
\end{align}
For an $H$ gate on an amplitude pair $(z_0,z_1)$, the new numerators are
$\sqrt2(z_0+z_1)$ and $\sqrt2(z_0-z_1)$. For $T$, multiply the bit-zero
numerator by two and the bit-one numerator by $\sqrt2(1+i)$.
For $S$, use two and $2i$ respectively. The $X$ and CNOT gates permute
amplitudes and multiply all numerators by two. Dividing by the new
$2^{L+1}$ gives precisely the usual gate actions.

The squared norm of one numerator is
\begin{equation}\label{eq:normquantum}
 |z|^2=A+B\sqrt2,\qquad
 A=a^2+2b^2+c^2+2d^2,\quad B=2(ab+cd).
\end{equation}
Thus for a selected set $J$ of measurement outcomes,
\begin{equation}\label{eq:probquantum}
 \Pr(J)=\frac{\sum_{j\in J}A_j+
                 (\sum_{j\in J}B_j)\sqrt2}{4^L}.
\end{equation}
The symbolic coefficient $B$ can be negative, although the final probability
is nonnegative.

\begin{proposition}[Decidable finite-circuit probabilities]\label{prop:finitequantum}
For a finite circuit over these gates with a computational-basis input,
comparison of a terminal basis-event probability with a rational threshold
is decidable by exact integer arithmetic. Its graph and its comparison
relations are Diophantine, with quartic representations.
\end{proposition}
\begin{proof}
The gate recurrences compute the four integer coordinates exactly.
After clearing positive denominators, a rational comparison reduces to the
sign of $A+B\sqrt2$ with $A,B\in\Z$. If either coefficient is zero or they
have the same sign, the sign is immediate. If they have opposite signs,
compare $A^2$ and $2B^2$, reversing the resulting sign when $A<0$.
Equality $A^2=2B^2$ with both coefficients nonzero is impossible: cancelling
common factors would give a reduced rational square root of two.
This supplies a terminating integer algorithm and therefore a computable
relation. Apply MRDP and Lemma~\ref{lem:quarticlower}.

For a direct bounded evaluator, use linear gate residuals and quadratic
norm residuals. Signed intermediate coordinates can be encoded by their
natural positive and negative parts; adding $z^+z^-=0$ makes that split
canonical. The finite comparison algorithm can be lowered to the same
arithmetic-gate format. This last observation does not assert the particular
witness count of Theorem~\ref{thm:singlefold} for the quantum compiler.
\end{proof}

\begin{example}[Cancellation and an irrational output probability]
Starting in $|0\rangle$, the circuit $H^2$ returns outcome zero with
probability one. If its intermediate basis alternatives were treated as
independent classical alternatives without recombination, the cancellation
in the outcome-one amplitude would be lost.

The circuit $HTH$, acting right to left on the input as usual, gives
\[
 \Pr(0)=\frac{2+\sqrt2}{4},\qquad
 \Pr(1)=\frac{2-\sqrt2}{4}.
\]
The accompanying sequential evaluator applies $H$, then $T$, then $H$, which
is the same palindromic circuit. Both probabilities are represented exactly,
not by decimal approximations. A finite fair-coin decision tree with a
uniform depth bound has only dyadic output probabilities, so no such tree
can realize this law exactly. The unbounded sampler in
Theorem~\ref{thm:semimeasure} can realize it and, because the two masses sum
to one, returns almost surely.
\end{example}

\subsection{Unbounded measured programs}
Consider an unbounded classical controller with finitely described quantum
gates and projective measurements, allowing only finitely many qubits at
each finite execution stage. A finite measurement history has a weight
computed by evolving an \emph{unnormalized} state through its gates and
measurement projections. Its squared norm is the history probability.
Keeping the state unnormalized avoids division by a possibly zero branch
probability.

Distinct recorded measurement histories are mutually exclusive classical
events. Summing the weights of first-return histories therefore gives a
presentation in Definition~\ref{def:firstreturn}. Lemma~\ref{lem:leftce} and
Theorem~\ref{thm:universal} apply. The same argument works for uniformly
computable gate entries when valid evolution and an effective finite
history presentation are part of the model assumptions; exact algebraic
zero tests are unnecessary for obtaining lower bounds.

\begin{corollary}[Quantum output-law Diophantine representation]
Every such effective quantum program's discrete classical terminal output
law has the fixed-quartic projected-root representation
\eqref{eq:rootmass}. Its ordinary strict output-probability thresholds are
existential Diophantine. Its conditional eventual-return thresholds are
subject to the distinctions in Section~\ref{sec:hierarchy}.
\end{corollary}
\begin{proof}
The measurement-history presentation gives a left-c.e. law; apply the
corresponding representation and threshold theorems. Classical fair-coin
control is a special case, so the general hardness obstructions cannot be
removed by admitting quantum operations.
\end{proof}

This is an exact statement about a terminal \emph{classical probability
law}. It is not an efficient classical simulation of quantum computation,
not a replacement for quantum states in interactive protocols, and not a
claim that all quantum behavior is captured by classical random variables.
The reference state-vector implementation uses $2^n$ amplitudes for $n$
qubits and is deliberately a small-instance exact checker. It does not
implement the unbounded universal MRDP compiler.

\section{An operator-valued normalization budget}\label{sec:operator}

\subsection{Accumulated ensembles}
A quantum analogue of coordinatewise accumulated output mass is a sequence
of positive semidefinite operators $\rho_s$, with
$\rho_s-\rho_{s-1}\succeq0$ and $\Tr\rho_s\le1$. For example, these may
represent sums of subnormalized density operators over disjoint returned
histories. Write
\[
 \rho_s=h_s\sigma_s,\qquad \Tr\sigma_s=1.
\]
This is an accumulation process, not the state trajectory of a single
system undergoing arbitrary unitary evolution.

For positive definite density matrices define
\begin{equation}\label{eq:K}
 K(\sigma,\tau)=\inf\{r>0:\sigma\preceq r\tau\}
 =\lambda_{\max}(\tau^{-1/2}\sigma\tau^{-1/2}).
\end{equation}
The logarithm of this quantity is the established max-relative entropy
$\Dmax(\sigma\Vert\tau)$, with logarithm-base conventions adjusted to our
natural logarithm \cite{datta}. The equality in \eqref{eq:K} follows by
conjugating the operator inequality by $\tau^{-1/2}$.

\begin{theorem}[Exact operator accumulation cost]\label{thm:operator}
For a finite or infinite sequence of positive definite density matrices
$\sigma_0,\sigma_1,\ldots$, put
\[
 K_s=K(\sigma_{s-1},\sigma_s),\qquad H_s=\prod_{j=1}^s K_j.
\]
There are accumulated ensembles $\rho_s=h_s\sigma_s$ with
$h_0=\delta>0$, $h_s\le1$, and $\rho_s-\rho_{s-1}\succeq0$ if and only if
$\delta H_s\le1$ for every $s$. Their unique pointwise minimal total masses
are $h_s^*=\delta H_s$. Equivalently, the logarithmic cost is
\begin{equation}
 \log H_N=\sum_{s=1}^{N}\Dmax(\sigma_{s-1}\Vert\sigma_s).
\end{equation}
\end{theorem}
\begin{proof}
A single monotonicity condition is
$h_s\sigma_s\succeq h_{s-1}\sigma_{s-1}$, equivalent to
$h_s/h_{s-1}\ge K_s$. Taking traces shows $K_s\ge1$.
The proof of Theorem~\ref{thm:budget} now applies verbatim to the scalar
factors: iteration gives necessity, and choosing equality at every step
gives positive semidefinite increments and sufficiency.
\end{proof}

For singular density matrices, the same result holds with $K_s=+\infty$
when $\supp(\sigma_{s-1})$ is not contained in
$\supp(\sigma_s)$, and with the inverse taken on the latter support
otherwise. To see the obstruction, a vector in the new kernel must also
lie in the old kernel whenever the old operator is bounded above by a
multiple of the new one. Under support inclusion, the finite-dimensional
restriction gives the same formula.

\begin{corollary}[Data processing and trace variation]\label{cor:operatorvariation}
For a completely positive trace-preserving map $\Lambda$,
\[
 K(\Lambda\sigma,\Lambda\tau)\le K(\sigma,\tau).
\]
Thus applying the same channel at each checkpoint cannot increase the
optimal accumulation cost. Moreover, any feasible operator path satisfies
\[
 \sum_{s=1}^{N}\frac12\|\sigma_s-\sigma_{s-1}\|_1
 \le\log(1/\delta).
\]
\end{corollary}
\begin{proof}
A positive map preserves $\sigma\preceq r\tau$, so every feasible $r$
before the map is feasible afterward. For the variation claim, if $K_s>1$
then
\[
 \sigma_s=K_s^{-1}\sigma_{s-1}+(1-K_s^{-1})\omega_s
\]
for a positive trace-one operator $\omega_s$. The triangle inequality and
$\|\omega_s\|_1=\|\sigma_{s-1}\|_1=1$ give trace distance at most
$1-K_s^{-1}$. Handle $K_s=1$ by equality, sum, and use the scalar logarithmic
inequality from Proposition~\ref{prop:variation}.
\end{proof}

The data-processing property of max-relative entropy is known
\cite{datta}. Here its immediate consequence is monotonicity of an entire
accumulation budget, rather than a new entropy monotonicity theorem.

\subsection{Two exact matrix examples}
\begin{example}[A coherent excursion erased by dephasing]\label{ex:coherent}
Let
\[
 \sigma_0=\sigma_2=I/2,\qquad
 \sigma_1=\begin{pmatrix}1/2&t\\t&1/2\end{pmatrix},\qquad0<t<1/2.
\]
The eigenvalues of $\sigma_1$ are $1/2\pm t$. Thus
\[
 K_1=\frac{1}{1-2t},\quad K_2=1+2t,\quad
 \delta_{\max}=\frac{1-2t}{1+2t}.
\]
At $t=1/6$ the cost is two and the maximum initial mass is $1/2$.
Dephasing in the displayed basis sends every matrix to $I/2$ and reduces
the cost to one. These matrices commute because the endpoints are scalar;
the example demonstrates a coherence-sensitive cost, not noncommutation.
\end{example}

\begin{example}[A genuinely noncommuting cycle]\label{ex:noncommuting}
Set
\[
 \rho=\begin{pmatrix}2/3&0\\0&1/3\end{pmatrix},\qquad
 \sigma=\begin{pmatrix}1/2&1/6\\1/6&1/2\end{pmatrix}.
\]
Their commutator is nonzero. Both have determinant $2/9$, and
\[
 \det(r\sigma-\rho)=\det(r\rho-\sigma)
 =\frac{4r^2-9r+4}{18}.
\]
The least admissible domination factor in either direction is
$K=(9+\sqrt{17})/8$. Indeed, the determinant vanishes there and the
Hermitian difference has positive trace, hence is positive semidefinite;
the smaller root is below one and cannot dominate a trace-one operator.
For the cycle $\rho\to\sigma\to\rho$,
\[
 G=K^2=\frac{49+9\sqrt{17}}{32},\qquad
 \delta_{\max}=\frac{49-9\sqrt{17}}{32}.
\]
Dephasing turns the middle state into $I/2$ and gives the strictly smaller
cycle cost $(4/3)(3/2)=2$. The symbolic checker verifies both determinant
identities and the reciprocal cost exactly.
\end{example}

The operator theorem is a finite-dimensional order optimization theorem.
It does not by itself compile every prescribed computable matrix sequence
into a circuit over the discrete gate set of Section~\ref{sec:quantum}.
Exact state preparation and a small Diophantine certificate for generalized
eigenvalue maxima are separate problems.

\section{Other unconventional substrates}\label{sec:substrates}

The decisive property for the output-law theorem is not a conventional
instruction set. It is an effective, nonovercounting presentation of finite
terminal histories with computable weights. Universality is needed for the
completeness lower bounds, but not for the positive representation theorem.
The following applications state their semantics explicitly.

\paragraph{Randomized combinator reduction.}
A lambda, SK, SKI, or Iota reducer with a specified effective reduction
strategy and a fair-coin branching primitive has finite syntactic
histories and dyadic history weights. Its first-return law therefore has
\eqref{eq:rootmass}. The random primitive and reduction strategy are part of
the proposed extension; this is not a claim that the inspected repository
already implements that probabilistic language. Adding probabilities does
not excuse ambiguities between deterministic evaluation strategy and
nondeterministic rewrite choice.

\paragraph{Stochastic fractional programs.}
One can give a FRACTRAN-style system an explicit rational random dispatcher
between finite deterministic fractional programs, preserving each program's
ordered divisibility semantics. A transition is then decidable arithmetic,
with a rational branch probability. Finite histories fit our interface.
This construction deliberately avoids saying ``choose an applicable
fraction randomly'' without specifying how conflicts, priorities, and
normalization are resolved. Deterministic FRACTRAN-to-Diophantine reductions
are already present in the literature \cite{larchey-forster}; the additional
object represented here is the probability law and its conditional boundary.

\paragraph{Random local rewriting and cellular systems.}
A finite description of each current active configuration, an effective
finite set of next local choices, and rational transition weights are enough
for an effective first-return presentation. Infinite configurations may
also be admitted when a finite-history simulation and observable can be
computed from their effective descriptions. Merely declaring an arbitrary
infinite initial configuration does not meet that requirement.

\paragraph{Continuous-time jump systems.}
For a countable-state jump process with finitely many outgoing rational
rates $q_{ij}$ at each effectively represented state, the next-jump
probability is $q_{ij}/\sum_kq_{ik}$ when the denominator is positive.
A zero-rate state is treated as absorbing. Events that reach a designated
terminal state after finitely many jumps therefore have an effective
rational-weight history presentation. This application uses the embedded
jump chain. It neither assigns a terminal value to an infinite explosion
history nor claims a result about arbitrary deadline-constrained events.
Those questions require additional semantics and analysis.

\paragraph{Robust observations of computable continuous trajectories.}
Suppose a trajectory is supplied with an algorithm for certified evaluation
at rational times and is continuous. Reaching an effectively open target
at some finite time is semidecidable: enumerate rational times and certified
interior balls, and accept once an evaluation is certified inside one.
Continuity ensures that an actual open-target visit contains an appropriate
rational time, including a one-sided interval at time zero. This yields a
Diophantine reachability predicate by MRDP. It says nothing about exact
zero crossings with no robust margin, or about differential equations whose
solutions are not uniformly computable from the supplied data.

\paragraph{The effective boundary is indispensable.}
An arbitrary real oracle parameter can encode a non-c.e. set, and an
observation that asks for an exact infinite-time limit can move beyond
existential arithmetic. Such a model is not covered merely because it is
called Turing-complete. The theorem concerns effective computational
substrates, not every imaginable extension that contains ordinary Turing
machines as a submodel.

\section{Reproducible checks and formalization route}\label{sec:audit}

\subsection{What the software actually verifies}
The bundle contains an exact rational budget evaluator, a symbolic
single-fold quartic compiler, a canonical prefix allocator, and an exact
small quantum-circuit evaluator. The script \code{code/verify.py} uses a
fixed random seed, $20260930$, for reproducible sampled tests. It also runs
exhaustive tests that do not rely on randomness.

\begin{center}
\small
\begin{tabular}{p{0.35\textwidth}p{0.55\textwidth}}
\toprule
Check & Observed result\\
\midrule
Exhaustive binary histories &
780 histories, 3,120 initial-mass cases, and 22,220 forced-pivot candidates;
exactly one valid pivot candidate precisely in each of the 1,330 feasible
cases.\\[3pt]
Multi-outcome and large-integer tests &
192 sampled schedules, including 32 additional rescalings by $10^{100}$;
all exact invariance and feasibility checks pass.\\[3pt]
Symbolic quartic &
$N=2,m=2$: 11 parameters, 33 witnesses, 38 residuals, degree four, 165
expanded monomials. All residuals vanish at the exported witness.\\[3pt]
Mutation checks &
Increasing each of the 33 witness coordinates separately by one is rejected.
This is a finite robustness check, not a search of all wrong tuples.\\[3pt]
Quadratic-field signs &
25,921 small integer pairs cross-checked with high-precision decimal
arithmetic; 200 near-cancelling Pell cases checked by exact identities.\\[3pt]
Quantum circuits &
100 sampled circuits of 30 gates; 3,000 exact norm checks pass, together with
$H^2$, $HTH$, and a Bell-state example.\\[3pt]
Prefix allocation &
12 dyadic approximation stages of the $HTH$ output law; disjoint allocated
mass is $4095/4096$.\\[3pt]
Operator budgets &
Four exact coherent examples and one noncommuting cycle pass symbolic
matrix checks.\\
\bottomrule
\end{tabular}
\end{center}

The two allocated masses at the twelfth $HTH$ stage are $437/512$ and
$599/4096$. Their sum is $4095/4096$, leaving exactly $1/4096$ unallocated
at that finite stage. No claim is made that twelve stages already represent
the full irrational law. The script also includes floating-point diagnostics
for the sharpness limit in Theorem~\ref{thm:binary}; those diagnostics are
explicitly labeled illustrative and do not support the exact proof.

The detailed results and software versions are in
\code{artifacts/verification.json}. The full finite quartic is supplied
both as a residual specification and as an expanded expression. The code
uses Python integers and \code{Fraction} for exact rational arithmetic;
SymPy is used for symbolic expansion and the matrix checks.

\subsection{Proof obligations suitable for Lean or Rocq}
A direct formalization should keep the finite arithmetic layer separate
from the analytic probability layer. The following dependency chain avoids
requiring a large quantum semantics development before obtaining a useful
kernel-checked result.

\paragraph{Finite order and ratio layer.}
Formalize positive rational rows, least maximizing indices, and the
single-step equivalence
$h'q_i\ge hp_i\ (\forall i)\iff h'\ge h\max_i(p_i/q_i)$.
Induct on the horizon to prove the optimal product theorem. This portion
requires finite sums, ordered fields, and finite maxima, not MRDP.

\paragraph{Certificate layer.}
Define a finite index type for parameters and witnesses, an explicit
residual family of multivariate integer polynomials, and its sum of
squares. Prove that natural zeros force every residual to vanish. Prove the
selector lemma, the product invariant, and witness uniqueness. Derive the
cardinality formulas from the index constructors rather than from an
external count. Prove a compiler correctness theorem and use the Python
example only as an untrusted certificate generator.

\paragraph{Computability and MRDP layer.}
Encode the allocation interpreter and the threshold semidecision procedure
on naturals. The inspected \code{Diophantine.mrdp} interface can then
produce the finite polynomial existence theorem. A separate circuit
lowering theorem establishes degree at most four. This layer should not
claim witness uniqueness for a generic MRDP construction.

\paragraph{Measure and limit layer.}
Prove the prefix-cylinder allocation invariant and its measure identity;
then identify limits of cumulative masses with output laws. Prove the
variation bounds and the weak-computability realization with its explicit
nonuniform existential tail choice. Separate program extraction claims from
existence in classical mathematics.

\paragraph{Quantum layer.}
First verify the integer-linear gate recurrences and the quadratic norm
formula, with the state-vector dimension explicit. Then develop measured
first-return semantics using unnormalized branch states. The
operator-budget theorem can be formalized independently in a finite
Hermitian matrix library.

This is a plan, not delivered Lean source or a report of successful Lean
compilation. In particular, no guessed theorem names or placeholder
\code{sorry} proofs are included as though they were checked results.

\subsection{A trust-boundary summary}
The ordinary proofs establish the general theorems. The finite tests provide
evidence about the supplied implementations, including arithmetic signs,
indexing, tie resolution, and degree accounting. They cannot replace an
inductive proof of the compiler theorem. Conversely, the MRDP existence
proof does not certify that the software has generated a universal
polynomial: it has not. The archive makes this separation explicit in both
the README and the machine-readable report.

\section{Further research questions}\label{sec:questions}

The following questions are proposed research directions, not claims that
their full priority or open-status history has been settled by the present
literature search. They are ordered to put concrete extensions ahead of
broader conjectural targets.

\begin{question}[Kernel-checked unique quartic compiler]
Can Theorem~\ref{thm:singlefold}, including invalid-input rejection and the
exact count $5Nm+m+4N+3$, be formalized and connected to an executable
polynomial exporter in ProveIt? A decisive deliverable would be a checked
reflection theorem accepting the supplied $N=2,m=2$ residual certificate.
\end{question}
The finite positivity and tie-breaking lemmas make this the most sharply
specified implementation target. It need not wait for quantum semantics or
infinite probability theory.

\begin{question}[Smaller canonical arithmetic certificates]
Can the $X,Y$ cross-product witnesses or the selectors be eliminated while
retaining degree at most four, natural-witness uniqueness, and a linear-size
arithmetic circuit? What trade-off between witness count, degree, and
coefficient size is optimal for this specific normalization predicate?
\end{question}
This is a finite compiler optimization problem. A mere MRDP existence
argument cannot answer its size or uniqueness requirements.

\begin{question}[Sharp multi-outcome variation budgets]
For a specified initial vector $p_0$, determine the sharp supremum of
$\sum_s\TV(p_{s-1},p_s)$ under
$\sum_s\Dinf(p_{s-1}\Vert p_s)\le B$. Does a useful potential formula
extend the balanced binary constant $1/2$ to $m\ge3$, and what paths are
asymptotically extremal?
\end{question}
Proposition~\ref{prop:variation} gives the universal upper bound $B$ but
Theorem~\ref{thm:binary} shows that it can be substantially nonsharp when
the starting law is fixed.

\begin{question}[Uniform high-success realization]
What effective additional data are necessary and sufficient to turn the
existential high-success realization in
Theorem~\ref{thm:weakclassification} into a uniform compiler? Natural input
formats include certified tail-variation bounds, an interior interval, or a
computable bound on the remaining logarithmic cost.
\end{question}
The present theorem explicitly leaves tail selection nonuniform. A solution
should specify its representation of real numbers and its promised inputs;
otherwise an apparent uniformization may hide an oracle.

\begin{question}[Exact symbolic infinite histories]
For which algorithmically described rational paths can $G_\infty$ be
computed, or bounded sharply enough to synthesize a desired success floor?
Examples include periodic transformations, linear recurrences, and
piecewise-rational paths with summable perturbations.
\end{question}
Finite prefix feasibility is decidable and uniquely certified here.
Certifying every prefix with one finite object is an additional mathematical
problem, not an automatic consequence of those certificates.

\begin{question}[Operator certificates with controlled multiplicity]
Can an explicit quartic natural-witness system certify the optimal factors
$K_s$ for algebraic density matrices, including support changes, without
uncontrolled multiplicity from eigenvector choices? Can canonical principal
minor or characteristic-polynomial data replace arbitrary spectral bases?
\end{question}
The scalar earliest-maximum rule has no immediate analogue for degenerate
eigenspaces. Any proposed uniqueness theorem must address that issue rather
than merely add spectral witnesses.

\begin{question}[Quantum front ends with structural compression]
Which circuit classes admit polynomial-size exact Diophantine certificates
that avoid the full $2^n$ state vector? Candidate interfaces include
stabilizer structure, bounded tensor-network width, and a bounded number of
non-Clifford gates, with all compression assumptions stated explicitly.
\end{question}
The current state-vector checker is exponential. A useful result would
preserve exact zero tests and Born probabilities while making the source of
any speedup mathematically visible.

\begin{question}[Interactive and compositional semantics]
How much of the projected-root representation survives for interactive
processes with adaptive inputs, correlated outputs, or sequential
conditioning? Can a compositional representation preserve joint laws rather
than only a terminal one-output marginal?
\end{question}
Independent samplers for marginals are not enough. A solution must preserve
the common probability space or the appropriate quantum instrument.

\begin{question}[Direct universal substrate polynomials]
Can one produce a fully expanded, audited quartic polynomial for a small
measured quantum controller or a stochastic fractional language, avoiding a
large generic interpreter? Which operation counts can be proved rather than
estimated from informal circuit descriptions?
\end{question}
Theorem~\ref{thm:universal} gives a fixed polynomial in principle, not its
coefficients or a competitive universal witness count.

\begin{question}[Counting executions without counting arithmetic witnesses]
Can canonical finite source certificates be composed into useful
resource-bounded counting reductions while preserving projected-event
weights? Under what explicit restrictions can one obtain a genuinely
parsimonious unbounded representation?
\end{question}
The finite single-fold construction here is not a general unbounded
single-fold theorem. Any extension must confront that distinction directly.

\begin{question}[Continuous-time and robust physical semantics]
Develop certified first-hit and deadline probabilities for effective jump
systems and computable flows, with explicit treatments of explosion,
Zeno behavior, exact equality observations, and robustness margins. Which
observables remain c.e. threshold events, and which cross the normalization
boundary proved here?
\end{question}
The finite-first-return interface identifies a tractable starting point;
it does not prescribe a physically preferred semantics for limit events.

\begin{question}[An algebra of probability specifications]
Build a specification language that tracks whether a quantity is left-c.e.,
right-c.e., weakly computable, or computable, and automatically assigns a
sound arithmetic quantifier class to each comparison. Can successful
normalization-budget certificates improve the inferred class for useful
restricted program families?
\end{question}
Such a system would prevent the common mistake of compiling every
probabilistic specification to a single existential equation merely
because its source language is Turing-complete.

\section{Conclusion}
The useful extension of Diophantine computation is not just another
unconventional universal instruction set. It is a semantics-preserving
account of which quantitative assertions admit finite arithmetic witnesses.
Effective output laws admit a common projected-root representation, but
normalization can raise a strict threshold from $\sigone$ to
$\sigtwo$, even with a success guarantee arbitrarily close to one.

The exact normalization budget supplies a constructive organizing theorem:
accumulated mass can follow a prescribed positive conditional path exactly
when its product of maximal coordinate ratios fits the available budget.
That theorem yields a direct, unique-witness quartic compiler, a sharp
binary oscillation bound, an exact weak-computability characterization,
and a noncommutative operator-order extension. The explicit finite compiler
and tests are ready to serve as a concrete formalization target; the broader
universal and quantum results retain their stated existence and semantic
boundaries.

\appendix
\section{Using the archive}\label{app:archive}
The archive is self-contained for reading the article and rebuilding its
PDF with a conventional LaTeX installation. Its main files are:
\begin{center}
\small
\begin{tabular}{p{0.42\textwidth}p{0.48\textwidth}}
\toprule
Path & Purpose\\
\midrule
\code{article.tex}, \code{article.pdf} & Source and typeset manuscript.\\
\code{code/normalization.py} & Exact budget evaluator and residual compiler.\\
\code{code/quantum_exact.py} & Exact ring arithmetic and small state vectors.\\
\code{code/prefix_allocator.py} & Canonical labeled prefix allocation.\\
\code{code/verify.py} & Reproducible test and artifact generator.\\
\code{artifacts/verification.json} & Test counts, versions, and limitations.\\
\code{README.md} & Build, execution, and scope instructions.\\
\bottomrule
\end{tabular}
\end{center}
The two files beginning \code{normalization_quartic_N2_m2} contain the
residual system with example assignment and the full expanded polynomial.
The file \code{HTH_prefix_allocation.json} records the finite prefix-code
example. Rebuild from the archive root with
\begin{verbatim}
python code/verify.py
pdflatex -interaction=nonstopmode -halt-on-error article.tex
pdflatex -interaction=nonstopmode -halt-on-error article.tex
\end{verbatim}
The Python scripts require Python 3.10 or newer and SymPy. Tests were run
with Python 3.13.5 and SymPy 1.14.0. A newly generated report records the
versions actually used. The source code uses zero-based coordinate indices;
the mathematical exposition uses indices $1,\ldots,m$. The earliest-maximum
rule is the same ordering in both conventions.

\section{Source and verification notes}
The repository was inspected at the pinned commit printed in
Section~\ref{sec:context}. Bibliography entries for repository sources use
that commit rather than a moving branch. The references on weakly computable
reals, probabilistic-program hardness, exact Clifford+$T$ arithmetic, and
max-relative entropy provide the established context; the proofs of the
specific budget and compiler statements appear in full in this manuscript.

No claim is made that the literature search establishes first publication
of every consequence proved here. In particular, general c.e.-to-Diophantine
existence, fair-coin representation of lower semicomputable output laws,
and the entropy formulas are not advertised as new discoveries. The primary
new deliverable of this work is the explicit synthesis of these ideas with
the canonical finite normalization certificate, including its verified
implementation and its carefully separated infinite semantics.

\begin{thebibliography}{99}
\small\raggedright
\setlength{\itemsep}{3pt}
\setlength{\parskip}{0pt}
\bibitem{proveit-mrdp}
Vladimir Reshetnikov and ProveIt contributors.
\emph{Diophantine/MRDP.lean}.
ProveIt repository, snapshot \snapshot, inspected 30 September 2026.
\href{https://github.com/VladimirReshetnikov/ProveIt/blob/f608f1cb3c5be8a736df1328c01b293aabfccf4e/Computability/HilbertTenthProblem/Lean/Diophantine/MRDP.lean}{Pinned source: \texttt{Diophantine/MRDP.lean}}.

\bibitem{proveit-readme}
Vladimir Reshetnikov and ProveIt contributors.
\emph{Hilbert's tenth problem: Jones's articles on Diophantine representation}.
Project README, pinned ProveIt snapshot, inspected 30 September 2026.
\href{https://github.com/VladimirReshetnikov/ProveIt/blob/f608f1cb3c5be8a736df1328c01b293aabfccf4e/Computability/HilbertTenthProblem/README.md}{Pinned project README}.

\bibitem{larchey-forster}
Dominique Larchey-Wendling and Yannick Forster.
Hilbert's Tenth Problem in Coq (Extended Version).
\emph{Logical Methods in Computer Science} 18(1), article 35, 2022.
\url{https://doi.org/10.46298/lmcs-18(1:35)2022}.
\url{https://arxiv.org/abs/2003.04604}.

\bibitem{ambos-spies}
Klaus Ambos-Spies, Klaus Weihrauch, and Xizhong Zheng.
Weakly computable real numbers.
\emph{Journal of Complexity} 16(4), 676--690, 2000.
\url{https://doi.org/10.1006/jcom.2000.0561}.
The bounded-variation characterization is also explicitly stated and
attributed in \cite{ng-stephan-wu}.

\bibitem{ng-stephan-wu}
Keng Meng Ng, Frank Stephan, and Guohua Wu.
Degrees of weakly computable reals.
In \emph{Logical Approaches to Computational Barriers}, CiE 2006.
\url{https://doi.org/10.1007/11780342_43}.
Author manuscript:
\url{https://personal.ntu.edu.sg/kmng/Files/Papers/ng-stephan-wu-cie-2006.pdf}.

\bibitem{kaminski-katoen}
Benjamin Lucien Kaminski and Joost-Pieter Katoen.
On the hardness of almost-sure termination.
MFCS 2015; author manuscript arXiv:1506.01930.
\url{https://arxiv.org/abs/1506.01930}.

\bibitem{giles-selinger}
Brett Giles and Peter Selinger.
Exact synthesis of multiqubit Clifford+$T$ circuits.
\emph{Physical Review A} 87, 032332, 2013.
\url{https://doi.org/10.1103/PhysRevA.87.032332}.
\url{https://arxiv.org/abs/1212.0506}.

\bibitem{datta}
Nilanjana Datta.
Min- and max-relative entropies and a new entanglement monotone.
\emph{IEEE Transactions on Information Theory} 55(6), 2816--2826, 2009.
\url{https://doi.org/10.1109/TIT.2009.2018325}.
\url{https://arxiv.org/abs/0803.2770}.
\bibitem{matiyasevich}
Yu. V. Matiyasevich.
Diophantine representation of enumerable predicates.
\emph{Mathematics of the USSR-Izvestiya} 5(1), 1--28, 1971.
\url{https://doi.org/10.1070/IM1971v005n01ABEH001004}.
\url{https://www.mathnet.ru/eng/im1910}.
\end{thebibliography}
\end{document}
